\documentclass[12pt,a4paper]{article}
\usepackage[utf8]{inputenc}
\usepackage[spanish]{babel}
\usepackage{amsmath,amssymb,amsfonts}
\usepackage{geometry}
\usepackage{hyperref}
\geometry{margin=2.5cm}

\title{\textbf{Constantes Físicas en Física Moderna}}
\author{Compendio de Constantes Fundamentales}
\date{\today}

\begin{document}

\maketitle

\section*{Introducción}
La física moderna comprende la mecánica cuántica, la física atómica, nuclear, de partículas y la relatividad. Sus constantes fundamentales gobiernan el comportamiento de la materia y la radiación a escala subatómica y astronómica.

\section{Constante de Planck ($h$ y $\hbar$)}

\subsection*{1. Significado, Símbolo y Explicación}
Simbolizada por $h$ (y $\hbar = \frac{h}{2\pi}$ para la constante reducida de Planck o constante de Dirac), es la constante de proporcionalidad fundamental que cuantiza la energía de los fotones ($E = h\nu = \hbar\omega$) y el momento angular orbital y de espín en mecánica cuántica. Es el pilar central del principio de incertidumbre de Heisenberg ($\Delta x \Delta p \ge \frac{\hbar}{2}$).

\subsection*{2. Valores Típicos en Ingeniería y Ciencia}
\begin{equation*}
h \approx 6.626 \times 10^{-34} \text{ J}\cdot\text{s} \quad \text{o} \quad 4.1357 \times 10^{-15} \text{ eV}\cdot\text{s}
\end{equation*}
\begin{equation*}
\hbar \approx 1.0546 \times 10^{-34} \text{ J}\cdot\text{s} \quad \text{o} \quad 6.5821 \times 10^{-16} \text{ eV}\cdot\text{s}
\end{equation*}

\subsection*{3. Decimales Completos y Definición Exacta (SI 2019)}
Desde el 20 de mayo de 2019, la constante de Planck se fijó exactamente para redefinir la unidad fundamental de masa (el kilogramo) mediante la balanza de Kibble:
\begin{equation*}
h = 6.62607015 \times 10^{-34} \text{ J}\cdot\text{s} \quad \text{(o kg}\cdot\text{m}^2\cdot\text{s}^{-1}\text{)}
\end{equation*}
\begin{equation*}
\hbar = \frac{6.62607015 \times 10^{-34}}{2\pi} \approx 1.0545718176461565 \times 10^{-34} \text{ J}\cdot\text{s}
\end{equation*}
\textbf{Incertidumbre / Margen de Error:} Exacto por definición ($u_r = 0$).

\vspace{0.8cm}
\hrule
\vspace{0.8cm}

\section{Masa del Electrón en Reposo ($m_e$)}

\subsection*{1. Significado, Símbolo y Explicación}
Simbolizada por $m_e$, es la masa estacionaria del electrón, la partícula elemental estable con carga negativa más ligera. Determina el radio de Bohr de los átomos, los niveles energéticos de Rydberg y la longitud de onda de Compton.

\subsection*{2. Valores Típicos en Ingeniería y Ciencia}
\begin{equation*}
m_e \approx 9.109 \times 10^{-31} \text{ kg} \quad \text{o} \quad 0.511 \text{ MeV}/c^2
\end{equation*}

\subsection*{3. Decimales Completos y Valor CODATA (2018/2022)}
\begin{equation*}
m_e = 9.1093837015(28) \times 10^{-31} \text{ kg}
\end{equation*}
En electronvoltios equivalentes:
\begin{equation*}
m_e c^2 = 0.51099895000(15) \text{ MeV}
\end{equation*}
\textbf{Incertidumbre / Margen de Error:}
\begin{align*}
\text{Incertidumbre absoluta: } & u(m_e) = 0.0000000028 \times 10^{-31} \text{ kg} \\
\text{Incertidumbre relativa: } & u_r = 3.0 \times 10^{-10}
\end{align*}

\vspace{0.8cm}
\hrule
\vspace{0.8cm}

\section{Constante de Estructura Fina ($\alpha$)}

\subsection*{1. Significado, Símbolo y Explicación}
Representada por la letra griega alfa ($\alpha$), es una constante física fundamental adimensional que caracteriza la fuerza de la interacción electromagnética entre partículas cargadas elementales:
\begin{equation*}
\alpha = \frac{e^2}{4\pi \varepsilon_0 \hbar c}
\end{equation*}
Determine la separación de la estructura fina en los espectros atómicos y es una constante de acoplamiento clave en electrodinámica cuántica (QED).

\subsection*{2. Valores Típicos en Ingeniería y Ciencia}
\begin{equation*}
\alpha \approx \frac{1}{137} \quad \text{o} \quad 0.00729735
\end{equation*}

\subsection*{3. Decimales Completos y Valor CODATA (2018/2022)}
\begin{equation*}
\alpha = 7.2973525693(11) \times 10^{-3}
\end{equation*}
Su inverso exacto ($\alpha^{-1}$):
\begin{equation*}
\alpha^{-1} = 137.035999084(21)
\end{equation*}
\textbf{Incertidumbre / Margen de Error:}
\begin{align*}
\text{Incertidumbre absoluta: } & u(\alpha) = 0.0000000011 \times 10^{-3} \\
\text{Incertidumbre relativa: } & u_r = 1.5 \times 10^{-10}
\end{align*}

\end{document}
